% GENERATED FILE. Edit canonical lessons, then run npm run generate:books.
% canonical-source-hash: fnv1a64:3825b897efbd5416
% canonical-lessons: MR-C119-shambhar, MR-C119-hajar, MR-C119-jag, MR-C119-rang, MR-C119-pravas

\chapter{Hundred, Thousand, World, Colour, Journey}
\label{ch:mr-hundred-thousand-world-colour-journey}
\hlchaptermodality{119}

\begin{chapteropening}
\textbf{By the end of this chapter:} \emph{I can say a hundred, a thousand, the world, a colour and a journey.}

Complete the last lesson of chapter 119: I can say a hundred, a thousand, the world, a colour and a journey.
\end{chapteropening}

\section[\texorpdfstring{\mr{शंभर} (shambhar)}{शंभर (shambhar)}]{\mr{शंभर} (shambhar) — a hundred}

\label{lesson:MR-C119-shambhar}



\begin{quote}
Before the new one: say the Marathi for seventy, then the Marathi for ninety.
\end{quote}

\subsection*{You'll want to know: \mr{शंभर}}
\textbf{\mr{शंभर}} — \emph{shambhar} — \textquotedblleft{}a hundred\textquotedblright{}.

\begin{cousinweb}
\textbf{\mr{शंभर}} (\emph{shambhar}) is a tadbhava, a word inherited and reshaped by sound change. It comes from \emph{*śatambhara}, \textquotedblleft{}full hundred\textquotedblright{}, a form CDIAL reconstructs from its descendants (CDIAL 12284). It is built on \emph{śata}, \textquotedblleft{}hundred\textquotedblright{}.
\end{cousinweb}

\begin{grammarlens}[title={What you've built}]
One more word for this chapter.
\end{grammarlens}

\subsection*{Your turn}
\begin{itemize}
\raggedright
  \item \emph{Say it:} \emph{shambhar}
  \item \emph{Say it:} \emph{shambhar}, once more
  \item \emph{Say it:} say \emph{shambhar}
  \item \emph{Recall:} say the Marathi for seventy, then the Marathi for ninety, then say \emph{shambhar} again
\end{itemize}

\begin{tcolorbox}[breakable,colback=teal!4,colframe=teal!35!black,title={Before you move on}]
What does \mr{शंभर} mean? (\textquotedblleft{}A hundred\textquotedblright{}.) Say it once more.
\end{tcolorbox}
\section[\texorpdfstring{\mr{हजार} (hajār)}{हजार (hajār)}]{\mr{हजार} (hajār) — a thousand}

\label{lesson:MR-C119-hajar}



\begin{quote}
Before the new one: say the Marathi for ninety, then the Marathi for a hundred.
\end{quote}

\subsection*{You'll want to know: \mr{हजार}}
\textbf{\mr{हजार}} — \emph{hajār} — \textquotedblleft{}a thousand\textquotedblright{}.

\begin{cousinweb}
\textbf{\mr{हजार}} (\emph{hajār}) is a loan, usually traced to Persian \emph{hazār}, \textquotedblleft{}thousand\textquotedblright{}. Marathi writes the \emph{z} as \textbf{\mr{ज}}.
\end{cousinweb}

\begin{grammarlens}[title={What you've built}]
One more word for this chapter.
\end{grammarlens}

\subsection*{Your turn}
\begin{itemize}
\raggedright
  \item \emph{Say it:} \emph{hajār}
  \item \emph{Say it:} \emph{hajār}, once more
  \item \emph{Say it:} say \emph{hajār}
  \item \emph{Recall:} say the Marathi for ninety, then the Marathi for a hundred, then say \emph{hajār} again
\end{itemize}

\begin{tcolorbox}[breakable,colback=teal!4,colframe=teal!35!black,title={Before you move on}]
What does \mr{हजार} mean? (\textquotedblleft{}A thousand\textquotedblright{}.) Say it once more.
\end{tcolorbox}
\section[\texorpdfstring{\mr{जग} (jag)}{जग (jag)}]{\mr{जग} (jag) — the world}

\label{lesson:MR-C119-jag}



\begin{quote}
Before the new one: say the Marathi for a hundred, then the Marathi for a thousand.
\end{quote}

\subsection*{You'll want to know: \mr{जग}}
\textbf{\mr{जग}} — \emph{jag} — \textquotedblleft{}the world\textquotedblright{}.

\begin{cousinweb}
\textbf{\mr{जग}} (\emph{jag}) is Sanskrit \emph{jagat}, \textquotedblleft{}moving, living\textquotedblright{}, which meant \textquotedblleft{}the world\textquotedblright{} from the Brahmanas on (CDIAL 5078). Marathi drops the final \emph{-t}.
\end{cousinweb}

\begin{grammarlens}[title={What you've built}]
One more word for this chapter.
\end{grammarlens}

\subsection*{Your turn}
\begin{itemize}
\raggedright
  \item \emph{Say it:} \emph{jag}
  \item \emph{Say it:} \emph{jag}, once more
  \item \emph{Say it:} say \emph{jag}
  \item \emph{Recall:} say the Marathi for a hundred, then the Marathi for a thousand, then say \emph{jag} again
\end{itemize}

\begin{tcolorbox}[breakable,colback=teal!4,colframe=teal!35!black,title={Before you move on}]
What does \mr{जग} mean? (\textquotedblleft{}The world\textquotedblright{}.) Say it once more.
\end{tcolorbox}
\section[\texorpdfstring{\mr{रंग} (raṅg)}{रंग (raṅg)}]{\mr{रंग} (raṅg) — a colour}

\label{lesson:MR-C119-rang}



\begin{quote}
Before the new one: say the Marathi for a thousand, then the Marathi for the world.
\end{quote}

\subsection*{You'll want to know: \mr{रंग}}
\textbf{\mr{रंग}} — \emph{raṅg} — \textquotedblleft{}a colour\textquotedblright{}.

\begin{cousinweb}
\textbf{\mr{रंग}} (\emph{raṅg}) may come from Sanskrit \emph{raṅga}, \textquotedblleft{}dye, colour\textquotedblright{} (CDIAL 10560), or from Persian \emph{rang}, \textquotedblleft{}colour\textquotedblright{}. The two words are alike in sound and sense and may have merged.
\end{cousinweb}

\begin{grammarlens}[title={What you've built}]
One more word for this chapter.
\end{grammarlens}

\subsection*{Your turn}
\begin{itemize}
\raggedright
  \item \emph{Say it:} \emph{raṅg}
  \item \emph{Say it:} \emph{raṅg}, once more
  \item \emph{Say it:} say \emph{raṅg}
  \item \emph{Recall:} say the Marathi for a thousand, then the Marathi for the world, then say \emph{raṅg} again
\end{itemize}

\begin{tcolorbox}[breakable,colback=teal!4,colframe=teal!35!black,title={Before you move on}]
What does \mr{रंग} mean? (\textquotedblleft{}A colour\textquotedblright{}.) Say it once more.
\end{tcolorbox}
\section[\texorpdfstring{\mr{प्रवास} (pravās)}{प्रवास (pravās)}]{\mr{प्रवास} (pravās) — a journey}

\label{lesson:MR-C119-pravas}



\begin{quote}
Before the new one: say the Marathi for the world, then the Marathi for a colour.
\end{quote}

\subsection*{You'll want to know: \mr{प्रवास}}
\textbf{\mr{प्रवास}} — \emph{pravās} — \textquotedblleft{}a journey\textquotedblright{}.

\begin{cousinweb}
\textbf{\mr{प्रवास}} (\emph{pravās}) is Sanskrit \emph{pravāsa}, \textquotedblleft{}dwelling abroad\textquotedblright{}, taken over whole as a learned loan, a tatsama (CDIAL 8802). It is \emph{pra-}, \textquotedblleft{}forth\textquotedblright{}, and \emph{vāsa}, \textquotedblleft{}dwelling\textquotedblright{}.
\end{cousinweb}

\begin{grammarlens}[title={What you've built}]
One more word for this chapter.
\end{grammarlens}

\subsection*{Your turn}
\begin{itemize}
\raggedright
  \item \emph{Say it:} \emph{pravās}
  \item \emph{Say it:} \emph{pravās}, once more
  \item \emph{Say it:} say \emph{pravās}
  \item \emph{Recall:} say the Marathi for the world, then the Marathi for a colour, then say \emph{pravās} again
\end{itemize}

\begin{tcolorbox}[breakable,colback=teal!4,colframe=teal!35!black,title={Before you move on}]
What does \mr{प्रवास} mean? (\textquotedblleft{}A journey\textquotedblright{}.) Say it once more.
\end{tcolorbox}
